\documentclass[a4paper,11pt]{article}
\usepackage[utf8]{inputenc}
\usepackage[T1]{fontenc}
\usepackage[english]{babel}
\usepackage{hyperref}
\usepackage{enumitem}
\title{KOTES401 — LLM‑avusteinen ohjelmoinnin harjoitus (workshop)}
\author{Opettaja}
\date{\today}
% CI rebuild: 2026-02-18
\begin{document}
\maketitle
\section*{Tavoite}
Opiskelija rakentaa pienen interaktiivisen komponentin (TODO‑lista) käyttäen kielimallia (LLM) apuna. Lopputulos on julkaisukelpoinen GitHub Pages\-sivuna.

\section*{Aikataulu (120 min)}
\begin{enumerate}[leftmargin=*]
  \item 0--10 min: Aloitus ja demo
  \item 10--20 min: Ympäristön avaus (Codespaces/GitHub.dev/paikallinen)
  \item 20--35 min: Prompt‑demo (koodin generointi, virheiden selvitys)
  \item 35--70 min: Päätehtävä (HTML/React‑komponentti)
  \item 70--90 min: Debuggaus ja refaktorointi LLM:llä
  \item 90--105 min: Lisäominaisuudet ja saavutettavuus
  \item 105--120 min: Demo‑kierros ja julkaisuohjeet
\end{enumerate}

\section*{Materiaalit}
\begin{itemize}
  \item Starter: \texttt{starter/} (HTML) ja \texttt{starter-react/} (React)
  \item Ratkaisu: \texttt{solution/} (opettajalle)
  \item Prompt‑mallit: \texttt{PROMPTS.md}
  \item Julkaisuohjeet: \texttt{PUBLISH.md}
\end{itemize}

\section*{Harjoitustehtävät}
Tavoite: rakentaa toimiva TODO‑komponentti (HTML tai React) ja osoittaa LLM:n hyödynnettävyys kehitysvaiheessa.

\subsection*{Perustehtävä}
\begin{enumerate}[leftmargin=*]
  \item Toteuta lis&auml;&auml;, poisto ja valmiiksi merkint&auml;.
  \item Tallenna lista \texttt{localStorage}iin siten, ett&auml; uudelleenlataus palauttaa tilan.
  \item Esittele v&auml;hint&auml;&auml;n yksi LLM‑interaktio (liitetty prompt + vastaus). 
  \item (J&auml;lkit&ouml;) Julkaise lopputulos kopioimalla tuotokset \texttt{public/}-kansioon.
\end{enumerate}

\subsection*{Hyv&auml;ksymiskriteerit}
\begin{itemize}
  \item Lis&auml;ys, poisto ja valmiiksi merkitseminen toimivat manuaalisesti.
  \item LocalStorage toimii (tiedot pysyvät sivun uudelleenlatauksen yli).
  \item V&auml;hint&auml;&auml;n yksi commit ja opiskelija selostaa LLM:n roolin muutoksessa.
\end{itemize}

\subsection*{Bonus‑teht&auml;v&auml;t}
\begin{itemize}
  \item Lisää yksikk&ouml;testi (esim. Jest / React Testing Library).
  \item Saavutettavuusparannuksia (ARIA, n&auml;pp&auml;imist&ouml;ohjaus).
  \item Lis&auml;&auml; muokkaus‑ tai suodatusominaisuus.
\end{itemize}

\subsection*{React‑osio (laajennettu)}
K&auml;yt&auml; \texttt{starter-react/}. Suositeltu komponenttirakenne: \texttt{TodoApp} (state), \texttt{TodoList}, \texttt{TodoItem}. Tilamuotoesimerkki:

\begin{verbatim}
[
  { id: 1, text: "Osta maitoa", done: false },
  { id: 2, text: "Soita äidille", done: true }
]
\end{verbatim}

Keskeiset funktiot: \texttt{addTodo(text)}, \texttt{removeTodo(id)}, \texttt{toggleTodo(id)}. Tallenna tila localStorageen esim.:
\begin{verbatim}
useEffect(() => {
  localStorage.setItem('todos', JSON.stringify(todos));
}, [todos]);
\end{verbatim}

Vinkki LLM‑avulle: pyyd&auml; LLM:lt&auml; event handlerien malli, yksikk&ouml;testit tai saavutettavuusehdotukset.

\section*{Esimerkkipromptit}
T&auml;ss&auml; muutamia valmiita prompt‑malleja opiskelijoille ja opettajalle — ja kaksi valmista esimerkkikooda + tulkinta -mallia, joita voi pyyt&auml;&auml; suoraan LLM:lt&auml;.
\begin{enumerate}[leftmargin=*]
  \item Kirjoita JavaScript‑funktio \texttt{addTodo(text)} joka lis&auml;&auml; uuden \texttt{li}‑elementin listaan \texttt{\#todo-list}. Lis&auml;&auml; my&ouml;s "Poista"‑painike joka poistaa kyseisen elementin. Palauta vain JavaScript‑koodi, ei selityksi&auml;.
  \item Korjaa virhe: "TypeError: Cannot read property \"dataset\" of null" — liit&auml; koodip&auml;tke&auml; ja pyyd&auml; lyhyt selitys ja korjaus.
  \item Refaktoroi funktio lyhyemmäksi ja kirjoita yksikk&ouml;testit (Jest): liit&auml; koodip&auml;tke&auml;.
  \item Kirjoita kuvaava git‑commit‑viesti muutokselle: lisätty Poista‑painike ja localStorage‑tuki.
  \item Paranna saavutettavuutta: "Lis&auml;&auml; ARIA‑attribuutit ja n&auml;pp&auml;imist&ouml;tuet t&auml;lle HTML‑pohjalle" (liit&auml; HTML‑p&auml;tke&auml;).
  \item Ymm&auml;rr&auml; koodi rivikohtaisesti: "Selit&auml; seuraava funktio rivi rivilt&auml; — kerro mahdolliset bugit ja parannusehdotukset" (liit&auml; koodip&auml;tke&auml;).
  \item Kirjoita testi ja CI‑workflow: "Lis&auml;&auml; Jest‑testit funktiolle X ja lis&auml;&auml; yksinkertainen GitHub Actions ‑workflow joka ajaa testit." 
\end{enumerate}

\subsection*{Esimerkkikoodit ja tulkinnat}
\paragraph{Esimerkki A — DOM‑poisto (bugi ja korjaus)}
\begin{verbatim}
function removeTodo(e) {
  const id = e.target.dataset.id;
  const item = document.querySelector(`#todo-${id}`);
  item.parentNode.removeChild(item);
}
\end{verbatim}
\textbf{Tulkinta:} Jos klikatun kohteen (e.target) rakenne vaihtelee (esim. sisällä on span tai svg), \texttt{e.target.dataset} voi olla undefined. My&ouml;s elementti voi olla jo poistettu — kutsu voi aiheuttaa virheen.
\textbf{Korjaus (esimerkki):}
\begin{verbatim}
function removeTodo(e) {
  const btn = e.target.closest('button');
  const id = btn?.dataset?.id;
  const item = document.querySelector(`#todo-${id}`);
  if (item) item.remove();
}
\end{verbatim}
\textbf{LLM‑prompt (esimerkki):} "Korjaa `removeTodo`-funktio siten, ett&auml; se toimii kun click‑target voi olla span tai svg napin sis&auml;ll&auml;. Palauta vain korjattu funktio."

\paragraph{Esimerkki B — React: tilan mutointi}
\begin{verbatim}
function toggleTodo(id) {
  const todo = todos.find(t => t.id === id);
  todo.done = !todo.done;
  setTodos(todos);
}
\end{verbatim}
\textbf{Tulkinta:} Koodi muuttaa `todos`-taulukon suoraan (mutaatio). React ei välttämättä havaitse muutosta, koska viite pysyy samana.
\textbf{Korjaus (esimerkki):}
\begin{verbatim}
function toggleTodo(id) {
  setTodos(todos.map(t => t.id === id ? { ...t, done: !t.done } : t));
}
\end{verbatim}
\textbf{LLM‑prompt (esimerkki):} "Korjaa `toggleTodo`-funktio: äl&auml; muuta tilaobjektia suoraan, palauta immuuttinen versio joka korjaa renderöinti‑ongelman."

\noindent Muokkausvihje: pyyd&auml; LLM:lt&auml; aina "show only the fixed code" tai "explain line-by-line" kun haluat koodin tulkinnan.

\section*{Julkaisu}
Julkaise kopioimalla tuotokset \texttt{public/}-kansioon ja tee commit.
\begin{enumerate}[leftmargin=*]
  \item Varmista, ett&auml; repo on \textbf{public} tai ett&auml; organisaatio sallii Pages‑julkaisun.
  \item Kopioi tiedostot (\texttt{index.html}, \texttt{styles.css}, \texttt{scripts/*.js}) \rightarrow \texttt{public/}.
  \item Commit & push:
    \begin{verbatim}
    git add public/*
    git commit -m "Publish: oma TODO app"
    git push origin main
    \end{verbatim}
  \item Odota Actions → "Deploy to GitHub Pages" ja tarkista Pages‑URL.
\end{enumerate}

\subsection*{PDF (workshop.tex -> workshop.pdf)}
Repossa on workflow PDF:n k&auml;ytt&ouml;&ouml;nottoon: \texttt{.github/workflows/latex-build.yml}. Triggerit: \texttt{push} ja \texttt{workflow_dispatch}. Artifactit: \texttt{workshop-pdf} ja \texttt{workshop-pdf-fallback}. Jos haluat, ett&auml; PDF commitoidaan automaattisesti repoon, anna write‑oikeudet Actionsille tai lisää secretti \texttt{PERSONAL_ACCESS_TOKEN} (scope: repo).

\section*{Opettajalle — vinkit ja aikataulu}
\begin{itemize}
  \item Demo &amp; prompt‑näyttö: 0–10 min (näytä LLM korjaamassa virhe).
  \item Ohjeistus &amp; työskentely: 10–70 min (perustehtävä + React‑ominaisuudet).
  \item Debuggaus &amp; refaktorointi LLM:llä: 70–105 min.
  \item Julkaisu &amp; demo: 105–120 min.
\end{itemize}

\section*{Liite: Vianm&auml;&auml;ritys ja CI}
\begin{itemize}
  \item "No jobs were run": tee uusi commit tai kaynnista workflow manuaalisesti (workflow_dispatch).
  \item Pages ei p&auml;ivity: varmista, ett&auml; muutokset menivt \texttt{public/}-kansioon ja Actions ajoi onnistuneesti.
  \item PDF ei ilmesty repoon: tarkista Actions‑artifactit; ota workflowlle write‑oikeudet tai lisaa \texttt{PERSONAL_ACCESS_TOKEN} -secretti.
  \item Workflow‑sijainti: \texttt{.github/workflows/latex-build.yml}.
\end{itemize}

\end{document}